\begingroup
\renewcommand{\arraystretch}{1.45}
\fancyhead[L]{\small HUNDERT SÄTZE ÜBER DIE NEUEN TRANSZENDENTEN.}

\Needspace{32\baselineskip}
\par\medskip\Needspace{5\baselineskip}\pdfbookmark[1]{Art. 4}{transscendentes-iv-art-4}\centerline{4.}\nobreak\smallskip
Lässt man $n$ ins Unendliche wachsen, so gehen die Reihen und Produkte des zweiten und dritten Satzes in unendliche Reihen und Produkte über. In dieser Form lautet der \textsc{vierte Satz}
\[
\begin{gathered}
1+x(y+\frac1y)+x^4(yy+\frac1{yy})+x^9(y^3+\frac1{y^3})+\text{etc. in inf.}\\[8pt]
=(1+xy)(1+\frac xy)(1+x^3y)(1+\frac{x^3}y)\\[8pt]\quad{}(1+x^5y)(1+\frac{x^5}y)\mathinner{.\,.\,.}\\[8pt]
\quad\times(1-xx)(1-x^4)(1-x^6)(1-x^8)\mathinner{.\,.}
\end{gathered}
\]

\Needspace{19\baselineskip}
und \textsc{der fünfte}
\[
\begin{gathered}
x^{\frac14}(y^{\frac12}+y^{-\frac12})+x^{\frac94}(y^{\frac32}+y^{-\frac32})\\[8pt]\quad{}+x^{\frac{25}4}(y^{\frac52}+y^{-\frac52})+\text{etc. in inf.}\\[8pt]
=x^{\frac14}(y^{\frac12}+y^{-\frac12}).(1+xxy)(1+x^4y)(1+x^6y)\mathinner{.\,.}\\[8pt]
\quad\times(1+\frac{xx}y)\quad(1+\frac{x^4}y)\quad(1+\frac{x^6}y)\mathinner{.\,.}\\[8pt]
\quad\times(1-xx)\quad(1-x^4)\quad(1-x^6)\mathinner{.\,.}
\end{gathered}
\]

\SourceLeaf{474}


\Needspace{28\baselineskip}
\par\medskip\Needspace{5\baselineskip}\pdfbookmark[1]{Art. 5}{transscendentes-iv-art-5}\centerline{5.}\nobreak\smallskip
Die Funktionen, die durch den vierten und fünften Satz in unendliche Produkte entwickelt werden, sind von großer Wichtigkeit, und es wird gut sein, sie hier durch besondere Funktionszeichen zu bezeichnen. Wir schreiben daher
\[
P(x,y)=1+x(y+y^{-1})+x^4(yy+y^{-2})+x^9(y^3+y^{-3})+\text{etc.}
\]
\[
\begin{gathered}R(x,y)=x^{\frac14}(y^{\frac12}+y^{-\frac12})+x^{\frac94}(y^{\frac32}+y^{-\frac32})\\[8pt]\quad{}+x^{\frac{25}4}(y^{\frac52}+y^{-\frac52})+\text{etc.}\end{gathered}
\]
zugleich auch
\[
Q(x,y)=1-x(y+y^{-1})+x^4(yy+y^{-2})-x^9(y^3+y^{-3})+\text{etc.}
\]

\Needspace{25\baselineskip}
wobei also $Q(x,y)=P(-x,y)=P(x,-y)$ gilt. Der einfachste Wert, den man $y$ geben kann, ist $1$. Da die entsprechenden Werte unserer Funktion von besonders großer Wichtigkeit sind und häufig vorkommen werden, schreiben wir der Kürze wegen anstelle von $P(x,1),Q(x,1),R(x,1)$ einfach $Px,Qx,Rx$. Wir bemerken noch, dass ein Exponent, der sich nur auf das Argument einer Funktion bezieht, durch Klammern gekennzeichnet wird, wie bei $P(x^3)$. Ohne Klammern ist stets anzunehmen, dass er sich auf die Funktion bezieht, sodass $Px^3$ dasselbe bedeutet wie $(Px)^3$. Also
\[
\begin{gathered}
Px=1+2x+2x^4+\mathord{.}\\[8pt]
Qx=1-2x+2x^4-\mathord{.}\\[8pt]
Rx=2x^{\frac14}+2x^{\frac94}+2x^{\frac{25}4}+\mathord{.}
\end{gathered}
\]

\Needspace{12\baselineskip}
Schließlich wollen wir in dieser Abhandlung mit dem Funktionszeichen $F$ das unendliche Produkt ausdrücken
\[
Fx=(1-x)(1-xx)(1-x^3)(1-x^4)\mathinner{.\,.\,.}
\]

\Needspace{22\baselineskip}
In diesen Zeichen lauten die beiden letzten Sätze
\[
\begin{gathered}\text{4.}\qquad P(x,y)=(1+xy)(1+\frac xy)(1+x^3y)(1+\frac{x^3}y)\\[8pt]\quad{}(1+x^5y)(1+\frac{x^5}y)\mathinner{.\,.\,.} Fxx\end{gathered}
\]
\[
\begin{gathered}\text{5.}\qquad R(x,y)=x^{\frac14}(y^{\frac12}+y^{-\frac12})(1+xxy)(1+\frac{xx}y)\\[8pt]\quad{}(1+x^4y)(1+\frac{x^4}y)\mathinner{.\,.\,.} Fxx\end{gathered}
\]
Und so ist offenbar
\[
\begin{gathered}\text{6.}\qquad Q(x,y)=(1-xy)(1-\frac xy)(1-x^3y)(1-\frac{x^3}y)\\[8pt]\quad{}(1-x^5y)(1-\frac{x^5}y)\mathinner{.\,.\,.} Fxx\end{gathered}
\]





\SourceLeaf{475}


\Needspace{22\baselineskip}
ferner, indem man $y=1$ setzt
\[
\begin{gathered}
7.\quad Px=(1+x)^2(1-xx)(1+x^3)^2(1-x^4)\\[8pt]\quad{}(1+x^5)^2(1-x^6)\mathinner{.\,.\,.}\\[8pt]
8.\quad Qx=(1-x)^2(1-xx)(1-x^3)^2(1-x^4)\\[8pt]\quad{}(1-x^5)^2(1-x^6)\mathinner{.\,.\,.}\\[8pt]
9.\quad Rx=2x^{\frac14}(1+xx)^2(1-xx)\\[8pt]\quad{}(1+x^4)^2(1-x^4)(1+x^6)^2(1-x^6)\mathinner{.\,.}
\end{gathered}
\]

\Needspace{18\baselineskip}
Setzt man hier $1+x=\dfrac{1-xx}{1-x}$, $1+x^3=\dfrac{1-x^6}{1-x^3}$ usw. ein, so gehen diese Ausdrücke in die folgenden über
\[
\begin{gathered}
10.\quad Px=\frac{(Fxx)^5}{(Fx)^2(Fx^4)^2}\\[8pt]
11.\quad Qx=\frac{(Fx)^2}{Fxx}\\[8pt]
12.\quad Rx=2x^{\frac14}\frac{(Fx^4)^2}{Fxx}
\end{gathered}
\]

\Needspace{29\baselineskip}
hieraus ergibt sich ferner
\[
\begin{gathered}
13.\quad Px.Qx=(Qxx)^2\\[8pt]
14.\quad Px.Rx=\tfrac12(R\mathord{\surd}{x})^2\qquad\text{oder was dasselbe ist}
\end{gathered}
\]
\[
Pxx.Rxx=\tfrac12(Rx)^2,\qquad Rxx=\frac{(Rx)^2}{2Pxx}
\]
\[
Qxx(Rx)^2=4x^{\frac12}.(Fx^4)^3\qquad\text{also}
\]
\[
\begin{gathered}15.\quad Fx\\[8pt]\quad{}=\mathord{{}^{3}\!\surd}{\frac{Q(x^{\frac12})(Rx^{\frac14})^2}{4x^{\frac18}}}
\\[8pt]\quad{}=\mathord{{}^{3}\!\surd}{\frac{P(x^{\frac12})Q(x^{\frac12})R(x^{\frac12})}{2x^{\frac18}}}
\\[8pt]\quad{}=\mathord{{}^{3}\!\surd}{\frac{(Qx)^2R(x^{\frac12})}{2x^{\frac18}}}\end{gathered}
\]

\Needspace{16\baselineskip}
ferner
\[
\begin{gathered}
16.\quad Px+Qx=2P(x^4)\\[8pt]
17.\quad Px-Qx=2R(x^4)
\end{gathered}
\]
Also durch Multiplikation nach 14
\[
18.\quad (Px)^2-(Qx)^2=2(Rxx)^2
\]

\Needspace{14\baselineskip}
Bezeichnet ferner $i$ die imaginäre Größe $\mathord{\surd}{-1}$, so gilt
\[
\begin{gathered}
19.\quad Px+iQx=(1+i)Q(ix)\\[8pt]
20.\quad Px-iQx=(1-i)P(ix)
\end{gathered}
\]

\Needspace{12\baselineskip}
Also durch Multiplikation

\SourceLeaf{476}

\[
21.\quad (Px)^2+(Qx)^2=2(Pxx)^2
\]
und aus der Multiplikation von 18 und 21 unter Heranziehung von 14
\[
22.\quad (Px)^4-(Qx)^4=(Rx)^4
\]

\Needspace{29\baselineskip}
Man sieht also, dass $(Pxx)^2$ das arithmetische Mittel zwischen $(Px)^2$ und $(Qx)^2$ ist. Da nach 13. die Größe $(Qxx)^2$ das geometrische Mittel zwischen denselben Größen darstellt und da (\textit{Theor. attract. el. p.}) bei zwei Größenreihen
\[
\begin{array}{llll}
m,&m',&m'',&m'''\mathinner{.\,.\,.}\\[8pt]
n,&n',&n'',&n'''\mathinner{.\,.\,.}
\end{array}
\]
die so verbunden sind, dass $m^{(\lambda)}$ stets das arithmetische und $n^{(\lambda)}$ das geometrische Mittel zwischen $m^{(\lambda-1)}$ und $n^{(\lambda-1)}$ ist, die gemeinsame Grenze das arithmetisch-geometrische Mittel von $m$, $n$ oder von irgendeinem Paar zusammengehöriger Größen der beiden Reihen genannt wird, ergibt sich der höchst wichtige \textsc{Satz} (23):

\textit{Das arithmetisch-geometrische Mittel zwischen $(Px)^2$ und $(Qx)^2$ ist stets $=1$.}


\Needspace{17\baselineskip}
Nach dem, was wir an der angegebenen Stelle bewiesen haben, ist also auch
\[
24.\quad\text{\textit{das Integral}}\qquad
\int\frac{\mathrm{d}\varphi}{\mathord{\surd}{((Px)^4\cos\varphi^2+(Qx)^4\sin\varphi^2)}}
\]
\textit{von $\varphi=0$ bis $\varphi=2\pi$ erstreckt $=2\pi$ oder auch von $\varphi=0$ bis $\varphi=\tfrac12k\pi$, wenn $k$ irgendeine ganze Zahl bedeutet, $=\tfrac12k\pi$.}

\clearpage
\endgroup
